\section{Introduction}
\label{sec:intro}

Full-shape analyses of galaxy clustering need the covariance matrix of the measured power-spectrum multipoles. Two routes are available. Sample covariances from large suites of mocks are free of modelling assumptions about the covariance itself, but they carry noise that propagates into parameter errors \cite{Hartlap2007,Taylor2013,Dodelson2013,Percival2014}, and they inherit every approximation of the mocks. Analytic covariances are noise-free and fast to recompute for a new footprint, weighting or cosmology, but they are only as good as the model. In practice the two are complementary: the analytic model is validated against mocks, and the mocks are used to check the model rather than to define the covariance.

This paper is about that validation step for the DESI \cite{DESI2016} DR2 power-spectrum multipoles, and about the methodology of the validation itself. The model is implemented in the public code \thecov\ \cite{thecov}. Its Gaussian term is computed from the survey randoms, with the full window function \cite{Grieb2016}. The non-Gaussian terms are super-sample covariance \cite{TakadaHu2013,Li2014}, including the local-average effect of normalising by the realised density, the two discreteness four-point terms, and, as a test, the connected trispectrum at tree level and in the response approximation. We compare it with 859 approximate (``holi'') mocks and 25 AbacusSummit $N$-body mocks \cite{Maksimova2021} of the first LRG redshift bin ($0.4<z<0.6$, LRG1) and of the quasars ($0.8<z<2.1$, QSO).

The paper is organised around two questions. The first: \emph{what does each term do?} Each term has a distinct structure in the space of the data vector, and these structures are what make the comparison with mocks informative. The Gaussian term is nearly diagonal. Super-sample covariance is essentially of rank two: coherent changes of the amplitude of the monopole and quadrupole. The discreteness terms are broad and smooth. The connected trispectrum couples all scales. The second question: \emph{how do we know?} A comparison of variances and $\chi^2$ distributions is a weak test of a covariance matrix with $n=\val{nbins}$ elements per side and $N=\val{Nmocks}$ mocks. Most of the information is in a few directions, the directions must be found without overfitting the mocks, and Gaussian-like errors of the model must be separated from missing non-Gaussian terms. We describe the set of diagnostics that achieves this and calibrate them against the noise of the mocks.

\Cref{sec:model} summarises the covariance model, \cref{sec:mocks} the mocks and measurements, and \cref{sec:method} the validation methodology. \Cref{sec:results} presents the results term by term. \Cref{sec:limits} discusses what these mocks can and cannot validate; in particular, the approximate mocks lack the long radial modes that dominate super-sample covariance in a thin redshift shell. \Cref{sec:conclusions} concludes.
